\documentclass[10pt,a4paper]{amsart}
\usepackage{amsmath,amssymb,amsthm,graphicx}
\usepackage{verbatim}
\newenvironment{reference}{\comment}{\endcomment}
%\newenvironment{reference}{}{}
\newenvironment{slogan}{\comment}{\endcomment}
\newenvironment{history}{\comment}{\endcomment}

% For commutative diagrams we use Xy-pic
\usepackage[all]{xy}

% We use 2cell for 2-commutative diagrams.
\xyoption{2cell}
% Theorem environments.
%
\theoremstyle{plain}
\newtheorem{theorem}[subsection]{Theorem}
\newtheorem{proposition}[subsection]{Proposition}
\newtheorem{lemma}[subsection]{Lemma}

\theoremstyle{definition}
\newtheorem{definition}[subsection]{Definition}
\newtheorem{example}[subsection]{Example}
\newtheorem{exercise}[subsection]{Exercise}
\newtheorem{situation}[subsection]{Situation}

\theoremstyle{remark}
\newtheorem{remark}[subsection]{Remark}
\newtheorem{remarks}[subsection]{Remarks}

\numberwithin{equation}{subsection}

% Macros
%
\def\lim{\mathop{\mathrm{lim}}\nolimits}
\def\colim{\mathop{\mathrm{colim}}\nolimits}
\def\Spec{\mathop{\mathrm{Spec}}}
\def\Hom{\mathop{\mathrm{Hom}}\nolimits}
\def\Ext{\mathop{\mathrm{Ext}}\nolimits}
\def\SheafHom{\mathop{\mathcal{H}\!\mathit{om}}\nolimits}
\def\SheafExt{\mathop{\mathcal{E}\!\mathit{xt}}\nolimits}
\def\Sch{\mathit{Sch}}
\def\Mor{\mathop{\mathrm{Mor}}\nolimits}
\def\Ob{\mathop{\mathrm{Ob}}\nolimits}
\def\Sh{\mathop{\mathit{Sh}}\nolimits}
\def\NL{\mathop{N\!L}\nolimits}
\def\CH{\mathop{\mathrm{CH}}\nolimits}
\def\proetale{{pro\text{-}\acute{e}tale}}
\def\etale{{\acute{e}tale}}
\def\QCoh{\mathit{QCoh}}
\def\Ker{\mathop{\mathrm{Ker}}}
\def\Im{\mathop{\mathrm{Im}}}
\def\Coker{\mathop{\mathrm{Coker}}}
\def\Coim{\mathop{\mathrm{Coim}}}

% Boxtimes
%
\DeclareMathSymbol{\boxtimes}{\mathbin}{AMSa}{"02}

%
% Macros for moduli stacks/spaces
%
\def\QCohstack{\mathcal{QC}\!\mathit{oh}}
\def\Cohstack{\mathcal{C}\!\mathit{oh}}
\def\Spacesstack{\mathcal{S}\!\mathit{paces}}
\def\Quotfunctor{\mathrm{Quot}}
\def\Hilbfunctor{\mathrm{Hilb}}
\def\Curvesstack{\mathcal{C}\!\mathit{urves}}
\def\Polarizedstack{\mathcal{P}\!\mathit{olarized}}
\def\Complexesstack{\mathcal{C}\!\mathit{omplexes}}
% \Pic is the operator that assigns to X its picard group, usage \Pic(X)
% \Picardstack_{X/B} denotes the Picard stack of X over B
% \Picardfunctor_{X/B} denotes the Picard functor of X over B
\def\Pic{\mathop{\mathrm{Pic}}\nolimits}
\def\Picardstack{\mathcal{P}\!\mathit{ic}}
\def\Picardfunctor{\mathrm{Pic}}
\def\Deformationcategory{\mathcal{D}\!\mathit{ef}}
\usepackage{cleveref}
\crefname{originallemma}{original lemma}{original lemmas}
\Crefname{originallemma}{Original lemma}{Original lemmas}
\makeatletter
\newlabel{lemma-colimit}{{09YN}{source}{}{sourcepointer-1}{}}
\newlabel{lemma-colimit@cref}{{[originallemma][0][]09YN}{[1][0][]source}{}{}{}}
\makeatother
\crefname{originallemma}{original lemma}{original lemmas}
\Crefname{originallemma}{Original lemma}{Original lemmas}
\makeatletter
\newlabel{categories-lemma-compose-counits}{{0DV0}{source}{}{sourcepointer-2}{}}
\newlabel{categories-lemma-compose-counits@cref}{{[originallemma][0][]0DV0}{[1][0][]source}{}{}{}}
\makeatother
\title{A sheaf on a colimit of sites is the colimit of its pushforwards and pullbacks}
\author{The Stacks Project Authors}
\date{stacks-project:sites.tex:local-snapshot-20261012; GNU FDL 1.2 or later}
\begin{document}
\maketitle
\noindent\textit{Editorial note.} All delivered mathematical text is the author's own, copied from the pinned Stacks Project chapter file sites-cohomology.tex; no statement or proof step was rewritten, reordered or removed.\par\medskip
\noindent\textit{Editorial note.} The proof is the author's complete argument; every step the author writes out is inside the counted span and no closing sentence is a placeholder.\par\medskip
\noindent\textit{Editorial note.} Every labelled reference occurring inside the counted spans is carried in assembly.external\_references as a non-load-bearing navigation pointer (same-chapter ones at their definition site, cross-chapter ones at the invocation inside the counted span), so no load-bearing mathematics is replaced by a stub. No bibliography citation occurs inside the counted spans.\par\medskip
\noindent\textit{Editorial note.} The checked chapter list for same-name collisions is the set of stacks chapters staged in this repository: chow, cohomology, divisors, etale-cohomology, flat, more-algebra, more-morphisms, morphisms, spaces-more-morphisms, varieties together with this card's chapter sites-cohomology.\par\medskip
\noindent\textit{Editorial note.} The preamble is the chapter's own retained minimal preamble (sources/stacks-sites-cohomology/retained-preamble.tex), a verbatim excerpt of the native preamble: the Xy-pic setup and the native theorem declarations with the mathematical macros.\par\medskip
\noindent\textit{Allowed standard prerequisites (editorial).}\begin{itemize}
\item Composition of counits of adjunctions. Categories, Lemma categories-lemma-compose-counits. Supplies the cocycle condition satisfied by the canonical maps.
\item The identification of a sheaf with the colimit of its pullbacks along the transition morphisms. Same chapter, Lemma lemma-colimit. Supplies the description of the values of the sheaf as a filtered colimit, which is compared with the colimit in the statement.
\end{itemize}\medskip
\ifdefined\hypertarget\hypertarget{sourcepointer-1}{}\fi
\noindent\textit{Source reference.} Original lemma 09YN: sources/stacks-sites-cohomology/pinned-original-source.tex, source lines 2905--2905. Navigation pointer invoked at proof line 4225; the counted span contains the pointer only and the cited result is not reproduced inside it, so the pointer is registered as a non-load-bearing navigation stub in the established form for Stacks material. No bibliography citation is involved.\par\medskip
\ifdefined\hypertarget\hypertarget{sourcepointer-2}{}\fi
\noindent\textit{Source reference.} Original lemma 0DV0: sources/stacks-sites/pinned-original-source.tex, source lines 4242--4242. Navigation pointer invoked at proof line 4242; the counted span contains the pointer only and the cited result is not reproduced inside it, so the pointer is registered as a non-load-bearing navigation stub in the established form for Stacks material. No bibliography citation is involved.\par\medskip
\begin{situation}
\label{situation-inverse-limit-sites}
Here we are given
\begin{enumerate}
\item a cofiltered index category $\mathcal{I}$,
\item for $i \in \Ob(\mathcal{I})$ a site $\mathcal{C}_i$ such that every
covering in $\mathcal{C}_i$ has a finite index set,
\item for a morphism $a : i \to j$ in $\mathcal{I}$ a morphism of sites
$f_a : \mathcal{C}_i \to \mathcal{C}_j$ given by a continuous functor
$u_a : \mathcal{C}_j \to \mathcal{C}_i$,
\end{enumerate}
such that $f_a \circ f_b = f_c$ whenever $c = a \circ b$ in $\mathcal{I}$.
\end{situation}

\begin{lemma}
\label{lemma-colimit-push-pull}
In Situation \ref{situation-inverse-limit-sites} assume
we have a sheaf $\mathcal{F}$ on $\mathcal{C}$. Then
$$
\mathcal{F} = \colim f_i^{-1}f_{i, *}\mathcal{F}
$$
where the transition maps are $f_j^{-1}\varphi_a$
for $a : j \to i$ where
$\varphi_a : f_a^{-1}f_{i, *}\mathcal{F} \to f_{j, *}\mathcal{F}$
is a canonical map
satisfying a cocycle condition as in Lemma \ref{lemma-colimit}.
\end{lemma}

\begin{proof}
For the morphism
$$
\varphi_a : f_a^{-1}f_{i, *}\mathcal{F} \to f_{j, *}\mathcal{F}
$$
we choose the adjoint to the identity map
$$
f_{i, *}\mathcal{F} \to f_{a, *}f_{j, *}\mathcal{F}
$$
Hence $\varphi_a$ is the counit for the adjunction given by
$(f_a^{-1}, f_{a, *})$. We must prove that for all
$a : j \to i$ and $b : k \to i$ with composition $c = a \circ b$
we have
$\varphi_c = \varphi_b \circ f_b^{-1}\varphi_a$.
This follows from Categories, Lemma \ref{categories-lemma-compose-counits}.
Lastly, we must prove that the map given by adjunction
$$
\colim_{i \in I} f_i^{-1}f_{i, *}\mathcal{F}
\longrightarrow
\mathcal{F}
$$
is an isomorphism. For an object $U$ of $\mathcal{C}$
we need to show the map
$$
(\colim_{i \in I} f_i^{-1}\mathcal{F}_i)(U) \to \mathcal{F}(U)
$$
is bijective. Choose an $i$ and an object $U_i$ of $\mathcal{C}_i$
with $u_i(U_i) = U$. Then the left hand side is equal to
$$
(\colim_{i \in I} f_i^{-1}\mathcal{F}_i)(U) =
\colim_{a : j \to i} f_{j, *}\mathcal{F}(u_a(U_i))
$$
by Lemma \ref{lemma-colimit}. Since $u_j(u_a(U_i)) = U$
we have $f_{j, *}\mathcal{F}(u_a(U_i)) = \mathcal{F}(U)$
for all $a : j \to i$ by definition. Hence the value of the colimit is
$\mathcal{F}(U)$ and the proof is complete.
\end{proof}

\end{document}
